% =====================================================================
% MÉTODOS – Datos: problemas de calidad del FASTA y cómo se resolvieron
% Persona A. Va dentro de sections/03-methods.tex (B coordina Métodos).
% =====================================================================

\subsection*{Calidad del archivo FASTA}

Antes de diseñar el \textit{pipeline} se exploró el archivo
\texttt{PanTEon\_Database\_v1.6.1.fasta} (947\,MB)~\cite{panteon2026zenodo}
con herramientas de línea de comandos que lo recorren línea por línea,
sin cargarlo en memoria. El archivo contiene 239\,946 secuencias de
2\,790 especies. Su estructura resultó consistente: todos los
encabezados siguen el patrón
\begin{center}\texttt{>ID\#CLASE/ORDEN/SUPERFAMILIA @Especie}\end{center}
\noindent todas las clasificaciones
tienen tres niveles, no hay identificadores repetidos, no hay secuencias
vacías (la más corta tiene 81\,pb) y el archivo usa saltos de línea Unix.
Aun así, se encontraron seis problemas, resumidos en la
Tabla~\ref{tab:calidad_fasta}.

\begin{table}[htbp]
\centering
\caption{Problemas de calidad del archivo FASTA y regla aplicada en el
\textit{pipeline}.}
\label{tab:calidad_fasta}
\small
\setlength{\tabcolsep}{4pt}
\begin{tabular}{p{3.7cm}p{2.6cm}p{4.9cm}}
\hline
\textbf{Problema} & \textbf{Alcance} & \textbf{Solución} \\
\hline
Secuencias en minúsculas & 208\,354 secuencias (86,8\,\%) &
Conversión a mayúsculas al leer, antes de calcular GC, N y
\textit{k}-mers. \\
Superfamilia escrita de dos formas (\texttt{SINE/TRNA} y
\texttt{SINE/tRNA}) & 56 secuencias &
Clasificación en mayúsculas; las superfamilias pasan de 66 a 65. \\
Nombres de especie con punto en las abreviaturas (\texttt{subsp.},
\texttt{sp.}), que en el CSV aparecen sin punto & 4 especies, 533
secuencias & Normalización común para FASTA y CSV (eliminar el punto,
espacios sobrantes y espacios múltiples). Con ella cruzan las 2\,790
especies. \\
Bases indeterminadas (N) & 54\,837 secuencias (22,9\,\%); 2\,147\,161 N
& Se conservan; se calcula \texttt{porcentaje\_n} y el \%GC se calcula
sobre A, C, G y T. No hay otros códigos IUPAC: el alfabeto válido es
ACGTN. \\
Identificadores que no siguen el formato \texttt{PDB\#\#\#\#\#\#\#\#}
& 31\,600 secuencias (13,2\,\%), de las especies de Dfam &
El campo \texttt{pdb\_id} guarda el identificador original sin asumir
un formato; el índice único se mantiene porque no hay repetidos. \\
Superfamilia igual al orden (\texttt{LINE/LINE}, \texttt{TIR/TIR}, etc.)
& 11 superfamilias, 21\,544 secuencias &
Se conservan tal cual y se documentan como superfamilia no determinada.
\\
\hline
\end{tabular}
\end{table}

Ningún registro se descartó: el \textit{pipeline} cuenta los registros
leídos, válidos y descartados por motivo (encabezado mal formado,
secuencia vacía o caracteres fuera de ACGTN), y en la ejecución sobre el
archivo completo los 239\,946 registros resultaron válidos. Se decidió
corregir en lugar de descartar porque cada problema afectaba a la forma
de los datos y no a su contenido: eliminar las secuencias con N, por
ejemplo, habría hecho perder casi una cuarta parte de la base. Además,
45 entradas del campo especie corresponden a taxones superiores (por
ejemplo, Homininae o Mammalia); se conservaron y se analizan en los
resultados como una limitación del modelo.

% =====================================================================
% MÉTODOS – Decisiones de diseño (pipeline del CSV y modelo de datos)
% =====================================================================

\subsection*{Decisiones de diseño}

\paragraph{Entidades y colecciones.}
Identificamos tres entidades, la especie, la secuencia y el perfil de
\textit{k}-mers, y cada una quedó en su propia colección. La colección de
especies tiene un documento por cada fila del CSV, con su taxonomía y el
origen de su biblioteca de TEs, mientras que las secuencias y los
\textit{k}-mers se guardan aparte porque son muchísimos más y tienen un
tamaño muy distinto. Decidimos separarlas así porque una especie tiene
muchas secuencias (hasta 1\,627 en una sola especie) y mezclarlas en un solo
documento habría hecho que cada consulta tuviera que cargar mucha
información que no necesita.

\paragraph{Información embebida y referenciada.}
La taxonomía de cada especie se guarda una sola vez en la colección de
especies y las secuencias solo guardan el identificador de su especie
(\texttt{especie\_id}), de modo que se llega a la taxonomía haciendo un
\texttt{\$lookup}. Esto porque embeber la taxonomía en cada secuencia habría
repetido los mismos datos cerca de 240\,000 veces, una por cada secuencia, y
cualquier corrección de un nombre habría que hacerla en todos los
documentos de esa especie en vez de uno. Otra opción era embeber toda la
taxonomía para evitar el \texttt{\$lookup}, pero la ganancia en las
consultas por reino no nos pareció suficiente para justificar esa
repetición. El costo de la referencia es que las consultas que agrupan por
reino o por filo necesitan unir las dos colecciones, y los identificadores
que genera MongoDB cambian cada vez que se recarga la colección de
especies, por lo que las secuencias también hay que volver a cargarlas.

\paragraph{Jerarquías de clasificación y de taxonomía.}
No se representa la jerarquía en los datos porque no creímos que fuera algo
que fuera importante almacenar, solo se pide que se pueda responder a
consultas y para ello no es necesaria una estructura jerárquica, así que se
dejó implícita. Los campos sin embargo están ordenados siguiendo esa
jerarquía, aunque esto no tenga relevancia al final más allá de un mejor
orden visual. Consideramos usar \textit{arrays} de modo que con el índice de
cada elemento se diera esta estructura, pero a la hora de hacer las
consultas es mucho más sencillo agrupar utilizando campos ya que no habría
que acceder a la posición.

\paragraph{Columnas del CSV y nombres de los campos.}
Solo se conservan las nueve columnas necesarias y las demás se descartan al
leer el archivo, ya que no aportan nada a las consultas ni al cruce. Los
nombres se pasaron al español y a \texttt{Class} y \texttt{Order} se les
puso el prefijo \texttt{tax\_} (\texttt{tax\_clase} y \texttt{tax\_orden})
porque la clase y el orden de la especie no tienen nada que ver con la clase
y el orden del elemento transponible que vienen en el FASTA, y dejarlos con
el mismo nombre habría sido una fuente de errores al hacer las consultas.

\paragraph{Limpieza de textos y comparación de nombres.}
A todas las columnas de texto se les quitan los espacios sobrantes al
inicio, al final y entre palabras, y los valores como \texttt{NA},
\texttt{N/A} o ``\texttt{-}'' se convierten en vacíos. Para comparar nombres
entre el CSV y el FASTA se usa una versión normalizada que pasa todo a
minúsculas y quita el punto de abreviaturas como \texttt{sp.} y
\texttt{subsp.}, pero esa versión solo sirve para comparar y nunca se
guarda, porque un nombre científico va con el género en mayúscula y al
exportar las bibliotecas el encabezado tiene que quedar igual que en el
archivo original. Otra opción era guardar el nombre normalizado, pero eso
habría hecho que las bibliotecas generadas ya no coincidieran con el FASTA.
Con esto el cruce funcionó para las 2\,790 especies, y solo 4 necesitaron la
normalización (una diferencia de un punto en \texttt{subsp.} o
\texttt{sp.}), que corresponden a 533 secuencias.

\paragraph{Valores faltantes, nombres inconsistentes y registros que no
cruzan.}
Para tratar los datos faltantes decidimos conservar las filas que tenían
campos vacíos excepto cuando falta el nombre de la especie, esto porque el
nombre de la especie es el principal identificador para relacionarla con las
secuencias del archivo FASTA y si este campo está vacío no es posible
realizar el cruce de forma confiable, por lo que en esos casos sí
descartamos la fila completa. Los demás atributos se conservan aunque no
tengan ningún valor, ya que no impide identificar la especie ni procesar su
secuencia.

Para los registros que no cruzan identificamos si los nombres normalizados
entre ambos archivos coinciden, de ser así, se le da prioridad al nombre del
archivo FASTA y se reemplaza el nombre del CSV. Decidimos hacerlo así porque
de igual manera habría que elegir cuál de los dos nombres debe utilizarse, y
modificar el CSV no solo es mucho más sencillo al haber menos lugares donde
realizar el cambio, sino que también si alguien pasara su propio archivo de
secuencias por nuestro programa para realizar alguna consulta este se vería
modificado y podría acarrear problemas si llegara a utilizarse en algún otro
lugar. Otra opción era simplemente ignorar estas discrepancias y no
contarlas en el análisis, pero un pequeño error en un nombre tiene mucho
menos peso que la pérdida de cientos de secuencias asociadas a ese nombre,
por lo que decidimos darle más importancia a esa información.

\paragraph{Duplicados.}
Comprobamos que el número de filas del CSV coincide con el número de
especies distintas después de normalizar (2\,790 en ambos casos), por lo que
no hay especies repetidas y no fue necesario fusionar filas. Decidimos no
escribir una lógica de fusión porque no habría tenido nada que fusionar y
solo habría añadido código que nadie usa, pero la comprobación se mantiene
en el reporte para que, si alguien cambia el CSV, el duplicado se note en
vez de perderse. Dejar que se elija la primera fila de cada grupo habría
sido más corto, pero habría descartado datos sin avisar.

\paragraph{Lectura del FASTA para el cruce.}
Se recorre el archivo línea por línea y solo se miran los encabezados, ya
que el FASTA pesa 947\,MB y cargarlo completo en memoria no era viable ni
necesario, porque para el cruce solo hace falta el nombre de la especie. El
costo es que el archivo se tiene que recorrer otra vez cuando se cargan las
secuencias, pero es pequeño comparado con el riesgo de quedarnos sin
memoria. Además, si una especie aparece escrita de dos formas distintas en
el FASTA el programa se detiene en vez de elegir una en silencio, porque esa
discrepancia no se puede resolver de forma automática y confiable.

\paragraph{Repetición del \textit{pipeline}.}
Antes de insertar se vacía la colección de especies, así que el
\textit{pipeline} se puede ejecutar varias veces y siempre deja el mismo
resultado. Otra opción era insertar solo lo nuevo, pero con 2\,790 filas es
mucho más simple borrar y volver a cargar todo, y eso evita que queden
documentos viejos si el CSV cambia.

\paragraph{Índices.}
En la colección de especies se usa un índice único sobre el nombre de la
especie y otro sobre la fuente. El primero impide que haya dos especies
iguales y agiliza las búsquedas por nombre que hacen varias consultas, y el
segundo sirve para la comparación por origen de los datos, donde se agrupa
por Dfam y RM2+MCH. El índice de \texttt{\_id} que MongoDB crea solo es el
que usa el \texttt{\$lookup} entre secuencias y especies. No creamos índices
sobre reino, filo o familia porque la colección tiene apenas 2\,790
documentos y recorrerla completa es casi instantáneo, así que un índice
extra solo ocuparía espacio sin que se notara la mejora.

\paragraph{Validación de esquema.}
Para la colección de especies no se utilizó una validación de esquema
(\texttt{\$jsonSchema}), esto principalmente por simplicidad: el
\textit{pipeline} ya hace validaciones que garantizan lo mismo que
garantizaría el esquema, así que no vimos necesaria esa capa extra de
protección y decidimos priorizar un código simple y legible.
